% 核心结论：电源接入完整串联回路，修复表头断线和Rx接线位置
% 草图：[真实几何与节点 | 条件、推导、结论]；17cm印宽，最小字号8pt。
% 图契约：欧姆表：闭合测量回路与非均匀刻度；图中的端口、几何、箭头和方程对应本题。
% 版面：沿用前三模块的几何图+推导分区；目标印宽17cm，简图至中等密度。
% 布局：左侧真实结构，右侧关键关系；真实电路边不作为架构流程箭头处理。
\documentclass[tikz,border=7pt]{standalone}
\usepackage{fontspec,xeCJK,amsmath,amssymb}
\IfFontExistsTF{Times New Roman}{\setmainfont{Times New Roman}}{\setmainfont{TeX Gyre Termes}}
\IfFontExistsTF{Heiti SC}{\setCJKmainfont[AutoFakeBold=1.5]{Heiti SC}}{\setCJKmainfont[AutoFakeBold=1.5]{FandolHei-Regular}}
\usetikzlibrary{arrows.meta,bending,calc,patterns,positioning,decorations.pathmorphing}
\definecolor{ink}{HTML}{243745}
\definecolor{blue}{HTML}{275DAB}
\definecolor{green}{HTML}{177C63}
\definecolor{red}{HTML}{BD4935}
\definecolor{orange}{HTML}{AA6C16}
\definecolor{purple}{HTML}{715393}
\tikzset{every node/.style={font=\fontsize{9.3}{12.5}\selectfont,text=ink},
arrow/.style={-{Stealth[length=5pt 1.5,width'=0pt 3]},shorten >=2pt,shorten <=1pt,line width=1pt},
arrow short/.style={-{Stealth[length=3pt 1.0,width'=0pt 2.5]},shorten >=1pt,shorten <=0pt,line width=.8pt},
axis/.style={arrow,draw=ink!60,line width=.6pt},
vec/.style={arrow,draw=blue},
wire/.style={draw=ink!75,line width=.8pt},
guide/.style={dashed,draw=ink!45,line width=.5pt},
card/.style={rounded corners=2pt,draw=blue!65,fill=blue!4,align=center,text width=6.5cm,inner xsep=9pt,inner ysep=7pt},
title/.style={font=\fontsize{10.5}{14}\selectfont\bfseries,text width=15cm,align=center},
label/.style={font=\fontsize{8.4}{11}\selectfont,align=center},
note/.style={label,text width=6.1cm}}
\begin{document}
\begin{tikzpicture}[x=1cm,y=1cm]
\path[use as bounding box] (-8,-3.5) rectangle(8,3.65);
\node[title] at (0,3.25) {\quad 欧姆表：闭合测量回路与非均匀刻度};
\begin{scope}[shift={(-4.05,0)}]
\draw[wire,fill=white](-1,1)circle(0.3);\node at(-1,1){$G$};
\draw[wire,fill=white](-2.3,0)circle(0.3);\node at(-2.3,0){};
\node at(-2.3,.14){$+$};\node at(-2.3,-.14){$-$};\node[left]at(-2.65,0){$\mathcal E$};\draw[wire](-2.3,.3)--(-2.3,1)--(-1.3,1)(-.7,1)--(-.2,1)(2.3,1)--(2.3,-1)--(1.2,-1)(-2.3,-.3)--(-2.3,-1)--(-1.2,-1);\draw[wire,draw=ink](-0.2,1)--(0.7749999999999998,1.0) (1.3249999999999997,1.0)--(2.3,1);\begin{scope}[shift={(1.0499999999999998,1.0)},rotate=0.0]\draw[draw=ink,fill=white,line width=.8pt](-0.275,-.085)rectangle(0.275,.085);\end{scope}
\node[label]at(1.05,1.42){$R_s$};
\draw[wire,draw=red](-1.2,-1)--(-0.275,-1.0) (0.275,-1.0)--(1.2,-1);\begin{scope}[shift={(0.0,-1.0)},rotate=0.0]\draw[draw=red,fill=white,line width=.8pt](-0.275,-.085)rectangle(0.275,.085);\end{scope}
\node[label]at(0.0,-0.58){$R_x$};
\node[label]at(0,-1.8){$R_0=R_g+R_s$；表头内阻已计入总内阻};\draw[ink](-2.2,-2.35)--(2.2,-2.35);\draw[ink](-2.2,-2.3)--(-2.2,-2.4);\node[label,above]at(-2.2,-2.35){$ 0$};\node[label,below]at(-2.2,-2.4){$ \infty$};\draw[ink](-1.1,-2.3)--(-1.1,-2.4);\node[label,above]at(-1.1,-2.35){$ 1/4$};\node[label,below]at(-1.1,-2.4){$ 3R_0$};\draw[ink](0,-2.3)--(0,-2.4);\node[label,above]at(0,-2.35){$ 1/2$};\node[label,below]at(0,-2.4){$ R_0$};\draw[ink](2.2,-2.3)--(2.2,-2.4);\node[label,above]at(2.2,-2.35){$ 1$};\node[label,below]at(2.2,-2.4){$ 0$};\end{scope}
\node[card,text width=6.5cm,draw=blue!65,fill=blue!4]at(3.9,1.1){$I=\mathcal E/(R_0+R_x)$};
\node[card,text width=6.5cm,draw=green!65,fill=green!4]at(3.9,-1.1){短接调零后：$I_g=\mathcal E/R_0$\\$R_x=R_0(I_g/I-1)$};

\end{tikzpicture}\end{document}